\section{Застосування, наслідки та обмеження}
\label{sec:case}

\subsection{Застосування та аналіз сценаріїв}

Інстанціюємо Задачу~\ref{prob:sizing} на $\approx3\times10^{7}$ бюлетенів за
дванадцять годин: середнє навантаження $\approx694$~TPS і, за припущеного
пік-фактора $\approx3$ (не виміряний профіль), пікове
$\Dpk=\resDemand$~TPS. Вимогу детекції --- скоординованої відмови з пропуском
транзакцій (coordinated transaction-omission fault) частки $\phi=0.2$ протягом
однієї епохи --- невиконувана, тож її послаблено (не застосовується;
підрозділ~\ref{sec:detection}), і $\nmin=4$. Прогони на CI без відмов і без
емульованої затримки фіксують не більше $\approx\resTpsMaxMeasured$~TPS (при
$\n=\resTpsMaxN$, $\cc=\resTpsMaxC$, $\q=\resTpsMaxQ$); модель працює в
одиницях \RNM{} $\TPSn=\TPS/\q$ (пропускна здатність на ядро квоти
валідатора), і за Припущенням~\ref{as:separable} масштаб $\kappaScale$ --- це
кількість ядер на валідатор у цільовому класі обладнання. Лінійність
пропускної здатності за квотою --- головне припущення й основний ризик
екстраполяції: $\kappaScale=\resKappa$ ядра у $\approx\resKappaQuotaRatio$
разів перевищує калібровану $\q=0.25$, тоді як X2 уже при $\q=0.33$ показує,
що пропускна здатність перестає зростати. Тож $\kappaScale$ --- вимога до
обладнання, яку треба перевірити на цільовому залізі, а не виміряна
можливість. За вільного
$\kappaScale$ Задача~\ref{prob:sizing} бере
\begin{equation}
\kappaScale^{\star}
=\frac{\Dpk}{\underline{\TPSn}_{1-\epsp}(4,\cc^{\star})},
\label{eq:kappa}
\end{equation}
мінімальний масштаб, при якому нижня межа прогнозного інтервалу при $\n=4$
досягає $\Dpk$.

\paragraph{Сценарій A (базовий).}
Підстановка каліброваних коефіцієнтів
($\beta=\resBeta$, $\gamma=\resGamma$) при $\cc^{\star}=\resCstar$ дає
$\kappaScale=\resKappa$ (ядер на валідатор) і вікно допустимості
$[\resNminCase,\resNmaxCase]$ з $\nopt=\resNoptCase$
(Пропозиція~\ref{prop:window}; рис.~\ref{fig:window}). Єдина допустима
кількість --- наслідок побудови: $\kappaScale$ --- найменший масштаб, що
допускає $\n=4$, тож $h(4)=0$ і, при $\beta>0$, $\nmax=4$; інформативним
результатом є саме $\kappaScale$. ($\cc^{\star}$ максимізує підігнану $\gfun$
поза сіткою МНК $\cc\le\resCalibCmax$; при $\cc=\resCalibCmax$
$\kappaScale=\resKappaCalibC$.)

\paragraph{Сценарії B і C (фіксований клас обладнання).}
Фіксуємо $\kappaScale=\resKappa$. Подвоєння $\Dpk$ дає $h(4)=-\log2<0$;
збільшення $\beta$ на WAN-приріст у межах X4 $\Delta\beta=\resDeltaBetaRtt$
знижує $h(4)$ на $\Delta\beta\log4$. Обидва вікна порожні
(Наслідок~\ref{cor:infeasible}), і додавання валідаторів не допомагає, бо
прогнозована пропускна здатність спадає з~$\n$. Обидва випливають з нульового
запасу в Сценарії~A; їхній зміст --- потрібний масштаб:
$\kappaScale=\resKappaDoubled$ і $\approx\resKappaWan$ для подвоєного попиту та
WAN відповідно.

Той самий робочий процес застосовується до реєстрів і ідентичності після
перекалібрування; обов'язкові вибори на блокчейні залишаються поза рекомендацією
\cite{springall2014,specter2020,park2021}.

\subsection{Наслідки}
\label{sec:discussion}

Підігнана пара $(\beta,\gamma)$ специфічна для Besu/QBFT і розгортання;
перенесення доречне лише тоді, коли шлях конкурентності, режим квот і клас RTT
збігаються з калібруванням, тоді як протокол \RNM{} і робочий процес
переносяться через перекалібрування. Бенчмарки мають звітувати $\lambda$ (або
$\lambda=0$), квоту на валідатор і клас RTT (Наслідок~\ref{cor:sign}).

\subsection{Обмеження}
\label{sec:limitations}

\emph{Клас відмов і детекція.} Зашифрований транспорт обмежує X5
omission/duplication, а отриманий детектор близький до випадкового; тож
детекційна частина задачі емпірично не задіяна, і демонстрація компромісу
«пропускна здатність --- детекція» потребує детектора з $J$, відділеним від
нуля.

\emph{Інфраструктура.} Чотириядерна CI з $\n\le\resNmaxMeasured$ ---
запланований субстрат; при $\n=\resNmaxMeasured$ резервація дорівнює чотирьом
ядрам хоста. Поза виміряними кількостями прогнози модель-базовані.
